\section{Tulkani --- ``The Itinerant Courts''}
\label{chap:tulkani}

\subsection*{Quick Reference}
\textbf{Tagline:} \textit{The knot remembers. The story pays. The circle holds.}

\begin{quote}
  \textbf{Caravan Judge (Elite):}  
  ``The road does not forgive a debt. It only reschedules it. A promise spoken under the open sky is witnessed by the wind, the grass, and the wheels of every cart that passes. I do not write judgments. I tie knots. The knots remember. When you are ready to pay, untie one. If you are not ready, tie another.''
\end{quote}

\begin{quote}
  \textbf{Wagon-Keeper (Commoner):}  
  ``My wagon is my hearth, my court, my grave. The Hollow cannot cross the circle of wheels unless invited. The tax collector cannot enter without a story. The lord cannot demand rent unless we camp on his land for nine nights. On the eighth night, we leave. That is the law.''
\end{quote}

\textbf{Genre / Mood:} Itinerant community, story-driven justice, folk horror at the edge of the Hollow.

\textbf{Starting Location:} A wagon circle at a crossroads, where nine painted wagons form a ring. A Caravan Judge sits on an overturned barrel, a knotted cord in her hand. ``You are standing on ground that belongs to the lord of the eastern valley. He will be here at dawn to demand rent. But you are in the circle. You are safe until the ninth hour. Tell me a story — a true one — and I will tell you how to leave before the lord arrives.''

\vspace{2mm}
\begin{tcolorbox}[colback=black!3,colframe=blue!60!black,title={Theme \& Atmosphere}]
The Tulkani have no fixed homeland; they are wanderers who follow the seasons and the trade routes. Their wagons are their hearths, their circles are their courts, and their knots are their ledgers. They are welcome everywhere and trusted nowhere — a position they wear with wry humour. A Tulkani caravan can appear at a lord's border, camp for seven nights, and leave before the eighth — because on the eighth night, the lord can demand rent. On the ninth night, the Hollow comes. The Tulkani never stay nine nights in one place. Their courts are circles of wagons; their judges are elders who carry braided cords. Each knot is a vow, a story, or a favour owed — but never merely a debt. They spread the cult of Ikasha — She Who Sleeps — not through temples, but through dream-cords tied in the dark. Their neutrality is sacred: a Tulkani camp is sanctuary for any traveller, regardless of allegiance, as long as they respect the circle and offer a story when asked.

\vspace{2mm}
\textbf{What you notice first:} The smell of woodsmoke and cooking oil. The painted wheels of wagons arranged in a circle, each wheel a different colour. The way everyone touches a knotted cord before speaking of an obligation.

\textbf{The one rule that kills newcomers:} Never sleep outside the circle at night. The wagon ring is sanctuary. Outside it, the Hollow walks. A Tulkani will not rescue you — they will only mourn you at dawn.
\end{tcolorbox}

\subsection*{Regional Mood: The Weight of the Knot}

Power here is measured in vow-cords, story‑credits, and the number of nights camped on a lord's land. The Tulkani have no king, no judges, no written law — only the Caravan Circle, where disputes are settled by storytelling. A liar is exposed not by evidence but by the knots on his cord: if a knot is tied in false witness, it will tighten until blood is drawn. Ikasha's cult is not a religion; it is a survival tactic — the dream inside the wagon, the whisper in the creaking wheel. The Tulkani pay no permanent taxes, but they pay tolls in stories. A lord who demands coin will be told a story so boring that his crops fail. A lord who offers hospitality will be told a story so beautiful that his harvest doubles. The Tulkani do not track debts; they track \textbf{reciprocity}. What you give, you get back — maybe tomorrow, maybe in the next valley. But the circle remembers.

\subsection*{Prominent NPCs of the Tulkani}
\label{sec:tulkani-npcs}

\begin{longtable}{|p{0.2\textwidth}|p{0.25\textwidth}|p{0.3\textwidth}|p{0.15\textwidth}|}
\hline
\textbf{Name} & \textbf{Role/Title} & \textbf{Motivation} & \textbf{Rumoured Quirk} \\
\hline
Elder Sarnai & Caravan Judge of the Nine Knots & Keep the circle whole & She ties a knot for every judgment; her cord drags on the ground. \\
\hline
Temur Wagon-Keeper & Master of the painted wheel & Find new grazing ground before the ninth night & He paints a new colour on his wagon each season; the colours are memories. \\
\hline
Gulshat & Story-Seller & Spread the dream of Ikasha through tales & She carries a dream-cord with nine knots; the ninth knot is loose. \\
\hline
The Knife (Unnamed) & Ikasha's agent & Free a spirit bound by the Veiled Sisters & He wears a mask made of knotted cord; he has no face underneath. \\
\hline
Vow-Broken & An exile who cut his own cord & Earn back his place in the circle & He carries a single knot tied in his own hair; the knot is his promise to return. \\
\hline
\end{longtable}

\subsection*{Names of the Tulkani}
\label{sec:tulkani-names}

Inspired by the nomadic traditions of the steppe and desert: short, guttural, with honorifics and clan prefixes. Surnames are often replaced by epithets describing a deed or a cord.

\begin{longtable}{|p{0.25\textwidth}|p{0.25\textwidth}|p{0.2\textwidth}|p{0.2\textwidth}|}
\hline
\textbf{First Names (Male)} & \textbf{First Names (Female)} & \textbf{Clan/Epithet} & \textbf{Knot Name} \\
\hline
Temur & Sarnai & \textit{the Wagon-Keeper} & \textit{Ninth Cord} \\
\hline
Batu & Gulshat & \textit{the Story-Teller} & \textit{Dream-Knot} \\
\hline
Orhan & Nergui & \textit{the Knife} & \textit{Unbound} \\
\hline
Cengiz & Altan & \textit{Vow-Broken} & \textit{Cut Cord} \\
\hline
Toghrul & Odval & \textit{Circle-Keeper} & \textit{Neutral Knot} \\
\hline
\end{longtable}

\paragraph*{(Circle/Wagon/Ground)} A wagon circle at dusk, lanterns lit; a painted wagon with a dream-cord hanging from the wheel; a crossroads where nine roads meet.

\section*{Spades --- Places}
\label{sec:tulkani-places}
\begin{enumerate}
\setcounter{enumi}{1}
\item \textbf{Wagon Circle (Night)} --- Nine wagons arranged in a ring, lanterns burning, wheels touching.
  \hfill \textit{Inside the circle, the Hollow cannot enter. Outside, it waits. The lanterns must never go out.} `[SANCTUARY]`
\item \textbf{Caravan Court (Day)} --- A carpet spread on grass; a judge on a barrel; a cord on the ground.
  \hfill \textit{The cord is the law. Step over it, and you accept the judge's verdict.} `[JUSTICE]`
\item \textbf{Dream-Wagon (Ikasha's Shrine)} --- A painted wagon with no wheels, sitting on stones. The door is a curtain of knots.
  \hfill \textit{Inside, a single candle and a bed of straw. Ikasha dreams there. Do not wake her.} `[MYSTIC]`
\item \textbf{Cord-Maker's Stall} --- A canvas awning where a woman braids horsehair and silk.
  \hfill \textit{Each cord is a promise. The colour tells the weight: white for truth, red for blood, black for a secret.} `[CRAFT]`
\item \textbf{Crossroads of the Ninth Road} --- A meeting of nine paths, only eight of which are visible.
  \hfill \textit{The ninth road is for the dead. The Tulkani leave an offering there each full moon.} `[THRESHOLD]`
\item \textbf{Abandoned Circle} --- A ring of burnt grass and wheel ruts. No wagons. No people.
  \hfill \textit{The circle was broken on the ninth night. The Hollow came. The grass will not grow back.} `[OMEN]`
\item \textbf{Vow-Cord Tree} --- An old oak with hundreds of knotted cords hanging from its branches.
  \hfill \textit{Each cord is a promise kept. The tree remembers. The wind unties one each season.} `[COMMUNITY]`
\item \textbf{Hollow's Gate (Seasonal)} --- A place where the boundary is thin, marked by a single painted stone.
  \hfill \textit{The Tulkani camp here only on the new moon. On the full moon, they are a day's ride away.} `[THRESHOLD]`
\item \textbf{Tinkanar (Itinerant Caravanserai)} --- A mobile market of wagons, open for one day each season.
  \hfill \textit{The location changes. The Tulkani announce it by dream — the dream of a sleeping merchant.} `[TRADE]`
\item[J] \textbf{Story-Well} --- A dry well where the Tulkani throw written stories on parchment.
  \hfill \textit{The stories are offerings to Ikasha. The well is full. The parchment never rots.} `[RITUAL]`
\item[Q] \textbf{The Ninth Cord's Post} --- A wooden post at the centre of every circle, with nine notches.
  \hfill \textit{The ninth notch is never carved. It is left for the Hollow. The Hollow never takes it.} `[SYMBOL]`
\item[K] \textbf{The Wandering Hearth} --- A wagon that is never hitched, always moving. It carries the council fire.
  \hfill \textit{The fire has been burning for three hundred years. It has never touched the ground.} `[LEGEND]`
\item[A] \textbf{Ikasha's Dreaming Place (Myth)} --- A cave beneath a hill, where She Who Sleeps lies on a bed of knotted cords.
  \hfill \textit{If you untie a cord, she wakes. If she wakes, the world ends.} `[MYTH]`
\end{enumerate}

\paragraph*{(Judge/Keeper/Storyteller)} Caravan Judge with a dragging cord; Wagon-Keeper painting a wheel; Story-Seller with a dream-cord and a cup of water.

\section*{Hearts --- People \& Factions}
\label{sec:tulkani-people}
\begin{enumerate}
\setcounter{enumi}{1}
\item \textbf{Caravan Judge} --- Old, grey-braided, a cord so long it drags on the ground.
  \hfill \textit{She ties a knot for every judgment. Untying a knot voids the judgment. She has never untied one.} `[JUSTICE]`
\item \textbf{Wagon-Keeper} --- A woman with paint on her hands, a brush behind her ear.
  \hfill \textit{She paints a new colour on her wagon each season. The colours are memories. She is running out of colours.} `[CRAFT]`
\item \textbf{Story-Seller (Gulshat)} --- A young woman with a dream-cord around her wrist. She speaks in rhymes.
  \hfill \textit{Her stories are true — but only if you listen with your eyes closed.} `[SOCIAL]` `[MYSTIC]`
\item \textbf{Ikasha's Knife} --- A hooded figure with a mask of knotted cord. No face underneath.
  \hfill \textit{He cuts the cords of oath-breakers. The cut cord becomes a ghost. The ghost follows the breaker.} `[JUSTICE]` `[DEATH]`
\item \textbf{Vow-Broken (Exile)} --- A man with a single knot tied in his own hair. He never smiles.
  \hfill \textit{He is saving stories to buy back his name. He needs nine. He has eight.} `[COMMUNITY]`
\item \textbf{Cord-Maker} --- A woman with braids of horsehair and silk between her fingers.
  \hfill \textit{She can tell the weight of a promise by the feel of the hair. Horsehair is heavy. Silk is light.} `[CRAFT]`
\item \textbf{Hollow-Watcher} --- A boy who sits at the edge of the circle at dusk, watching the dark.
  \hfill \textit{He can see the Hollow when it walks. He points. The adults throw salt.} `[SCOUT]`
\item \textbf{The Dreamer} --- A woman who sleeps in the Dream-Wagon. She never wakes during the day.
  \hfill \textit{She dreams the caravan's path. The wagon follows her dreams.} `[MYSTIC]`
\item \textbf{Tinkanar Merchant} --- A fat man with nine bells on his hat. He sells knots and stories.
  \hfill \textit{He accepts only promises as payment. His ledger is a cord. It has nine thousand knots.} `[TRADE]`
\item[J] \textbf{The Ghost of the Ninth Night} --- A translucent figure seen at abandoned circles.
  \hfill \textit{It was a Tulkani who stayed too long. It asks: ``Is it the ninth hour yet?''} `[DEATH]`
\item[Q] \textbf{The First Wagon-Keeper (Legend)} --- A skeleton driving a wagon of bones. The wheels are skulls.
  \hfill \textit{It appears when a circle is broken. It offers a ride. The ride leads to the Hollow.} `[LEGEND]`
\item[K] \textbf{Ikasha's Voice} --- A whisper heard in the Dream-Wagon at midnight. It speaks in riddles.
  \hfill \textit{``The ninth cord is not a cord,'' it says. ``It is the space between the knots.''} `[MYSTIC]`
\item[A] \textbf{The Unwitnessed} --- A Tulkani who never tied a single knot. No one remembers her name.
  \hfill \textit{She walks through circles unnoticed. She can break any vow because no one saw her make it.} `[MYTH]`
\end{enumerate}

\paragraph*{(Toll/Story/Hollow)} A lord demands rent for camping on his land — negotiate nights or tell a story; the Hollow walks on the ninth night — a character hears their own name whispered; a stranger recognises your vow-cord and asks for a tale.

\section*{Clubs --- Complications/Threats}
\label{sec:tulkani-complications}
\begin{enumerate}
\setcounter{enumi}{1}
\item \textbf{Recognition (Vow-Cord Seen)} --- Someone from your past sees the knots you carry. They ask for a story you owe. `[SOCIAL]` `[RECIPROCITY]`
  \hfill \textit{The story must be about a promise you kept — or broke. If you cannot tell one, they tie a new knot.}
\item \textbf{Lord's Rent Demand} --- The local lord arrives at the circle at dawn. ``You have been here seven nights. Pay the eighth.'' `[AUTHORITY]` `[NEGOTIATION]`
  \hfill \textit{The price is a story that flatters him. If the story fails, he seizes a wagon. The circle holds its breath.}
\item \textbf{Hollow Walks (Ninth Night)} --- Someone in the circle has miscounted. The ninth hour approaches. `[DEATH]` `[FEAR]`
  \hfill \textit{The lanterns flicker. A shadow stands at the edge. It whispers a name. The name is the last person who broke an oath.}
\item \textbf{Vow-Cord Tangled} --- Two promises conflict. The knots tighten. Blood is drawn. `[CONFLICT]` `[RECIPROCITY]`
  \hfill \textit{The judge must cut one cord. The severed cord becomes a ghost. The ghost accuses the cutter of favouritism.}
\item \textbf{Wagon Wheel Broken} --- A wheel cracks. The circle is incomplete. The Hollow can enter. `[DISASTER]`
  \hfill \textit{The wheel must be repaired before dusk. The repair requires a story — a true one — told to the wood. The wood listens.}
\item \textbf{Tulkani Schism} --- Two judges dispute the meaning of a knot. The circle splits. `[POLITICS]`
  \hfill \textit{Each judge claims authority. The party must choose which cord to follow. The other cord will remember — and so will the Hollow.}
\item \textbf{Ikasha's Dream Disturbed} --- A stranger enters the Dream-Wagon. The Dreamer wakes. `[MYSTIC]` `[CRISIS]`
  \hfill \textit{Ikasha stirs. The ground trembles. The stranger must tell a story so good that Ikasha falls back asleep. The stranger is you.}
\item \textbf{Story Toll (Bridge/Pass)} --- A guard demands a story to cross. Your story is not good enough. `[TRAVEL]` `[SOCIAL]`
  \hfill \textit{The guard demands a truth instead. The truth must be a secret. The secret must hurt. The guard will remember.}
\item \textbf{Vow-Broken's Return} --- An exile tries to rejoin the circle. The judge refuses. `[COMMUNITY]` `[DRAMA]`
  \hfill \textit{The exile demands the party speak for him. If they refuse, he cuts his own throat. If they speak, they owe a knot — and a story.}
\item[J] \textbf{The Ninth Cord Cut} --- Someone cuts the ninth knot on the central post. The Hollow floods in. `[DEATH]` `[TRAP]`
  \hfill \textit{The lanterns go out. The circle becomes a cage. The party must re-tie the knot in the dark. The cord is warm.}
\item[Q] \textbf{Ikasha's Knife Betrayed} --- The Knife is offered a bribe to cut a false cord. He refuses. The briber attacks. `[COMBAT]` `[JUSTICE]`
  \hfill \textit{The Knife is wounded. He asks the party to carry his mask. The mask weighs nothing — until you put it on. Then it weighs everything.}
\item[K] \textbf{Tinkanar Closed (No Stories)} --- The mobile market is empty. The merchant has run out of knots. `[TRADE]` `[CRISIS]`
  \hfill \textit{He is selling his own name. The price is a story that makes him weep. He has not wept in forty years. The last story he heard was a lie.}
\item[A] \textbf{Ikasha Wakes} --- The Dreamer dies. Ikasha opens her eyes. The sky turns white. `[APOCALYPSE]`
  \hfill \textit{She asks one question: ``Who cut the ninth cord?'' The answer must be a name. If no name is given, she takes all names.}
\end{enumerate}

\paragraph*{(Cord/Bead/Charm)} Vow-bead (a single knot, can be used to call in a favour once); safe-haven oath (a knotted cord that guarantees sanctuary in any Tulkani circle); dream-cord (gift from Ikasha, one question answered in a dream).

\section*{Diamonds --- Rewards/Leverage}
\label{sec:tulkani-rewards}
\begin{enumerate}
\setcounter{enumi}{1}
\item \textbf{Vow-Bead (Single Knot)} --- A small knot on a leather thong. Present it to call in a favour from any Tulkani. `[SOCIAL]`
  \hfill \textit{The knot will loosen after use. If it is not retied within a year, the favour becomes a story owed — and stories are harder to repay.}
\item \textbf{Safe-Haven Oath (Knotted Cord)} --- A braided cord with three knots. Hang it on your wagon to guarantee sanctuary in any Tulkani circle. `[TRAVEL]`
  \hfill \textit{The cord must be renewed each season. A new knot is added. After nine knots, it becomes a dream-cord — and the Hollow will notice you.}
\item \textbf{Dream-Cord (Ikasha's Gift)} --- A silk cord with nine knots, the ninth loose. Once, sleep on it to ask Ikasha one question. She will answer in a dream. `[MYSTIC]`
  \hfill \textit{The answer is always a riddle. The ninth knot tightens after use. When it tightens fully, Ikasha wakes — and the world changes.}
\item \textbf{Caravan Judge's Cord (Borrowed)} --- A length of braided horsehair. Tie it around your wrist to gain +1 Position in any negotiation with Tulkani. `[SOCIAL]`
  \hfill \textit{The cord is the judge's own. She will want it back. The return is a negotiation — and a story.}
\item \textbf{Wagon-Keeper's Paint (Pot)} --- A clay pot of blue pigment. Paint your own wagon wheel, and no lord will demand rent for nine days. `[TRAVEL]`
  \hfill \textit{The paint is made from crushed sky-stone. It glows under moonlight. The Hollow avoids it — but the lord will send a spy.}
\item \textbf{Story-Seller's Verse} --- A single verse of a true story. Recite it to calm a spirit or end a feud. `[SOCIAL]` `[MYSTIC]`
  \hfill \textit{The verse works once. After that, the story belongs to the listener. You cannot tell it again — but you can tell a different truth.}
\item \textbf{Ikasha's Mask (Knotted)} --- A mask of knotted cord. Wear it to become invisible to the Hollow for one night. `[STEALTH]` `[DEATH]`
  \hfill \textit{The mask has no eye holes. You must walk by memory. The memory may fail. If it fails, the Hollow sees you — and smiles.}
\item \textbf{Hollow-Watcher's Salt Pouch} --- A leather pouch of salt blessed by the boy who watches the circle. Throw it to banish a Debt-Walker. `[COMBAT]` `[DEATH]`
  \hfill \textit{The salt is bitter. The Debt-Walker will taste it. It will remember your face — and your smell.}
\item \textbf{Tinkanar Merchant's Bell} --- One of nine bells from the merchant's hat. Ring it to summon him anywhere. `[TRADE]`
  \hfill \textit{He will come. He will demand payment in stories. He will not leave until paid. He has a long memory.}
\item[J] \textbf{Vow-Broken's Hair Knot} --- A single knot tied in human hair. It contains a broken promise. Use it to void a contract. `[SOCIAL]` `[MAGIC]`
  \hfill \textit{The knot is from an exile. He wants it back. He will pay a story for it — the story of how he broke his vow.}
\item[Q] \textbf{The Ninth Cord's Post Fragment} --- A splinter of wood from the central post. Touch it to any knot to learn who tied it. `[MYSTIC]`
  \hfill \textit{The splinter is warm. It will burn your hand if the knot was tied in false witness. The truth leaves no scar.}
\item[K] \textbf{First Wagon-Keeper's Skull} --- A human skull painted with wheel patterns. Place it at the centre of a circle to keep the Hollow away for nine nights. `[DEATH]` `[SANCTUARY]`
  \hfill \textit{The skull whispers on the ninth night. It asks: ``Is it time?'' Say no. If you say yes, it will take your place.}
\item[A] \textbf{Ikasha's Dream (Bottle)} --- A glass bottle containing a sleeping breath. Break it to send a question to Ikasha. She will answer in a vision. `[MYSTIC]`
  \hfill \textit{The vision lasts one minute. The cost is one memory. The memory is forgotten forever. The memory is always a happy one.}
\end{enumerate}

\section*{Quick use notes}
\label{sec:tulkani-quick-use}
\begin{itemize}
\item Draw until you have all four suits: Spade = place, Heart = actor, Club = pressure, Diamond = leverage. Highest rank sets the main Clock (2–5 → 4, 6–10 → 6, J/Q/K → 8, A → 10).
\item Diamonds are codified outcomes (knots, cords, beads) that change position rather than call for a roll.
\item If any A appears, echo \textbf{knot-omens}—a cord that tightens without being pulled, a knot that unties itself, a wheel that stops turning at the ninth spoke.
\end{itemize}

\subsection*{Additional Features: Itinerant Mechanics}
\label{sec:tulkani-itinerant}

\subsubsection*{The Ninth Night Rule}
A Tulkani caravan may camp on a lord's land for seven nights without payment. On the eighth night, the lord may demand rent — typically a story, a promise, or a small service. On the ninth night, the Hollow comes. No Tulkani stays nine nights in one place. A party camping with Tulkani must leave by the eighth dusk or face the Hollow's attention.

\textbf{Mechanic:} Each night the party camps in a Tulkani circle, mark a segment on a \textbf{Ninth Night Clock [8]}. At 7 segments, the lord (or local authority) arrives to demand rent. At 8 segments, the Hollow manifests. The party can leave before 8 to reset the clock.

\subsubsection*{Story as Currency}
Tulkani do not use coin. They trade in stories, promises, and knots. A good story can buy a meal, a horse, or safe passage. A bad story can start a feud. The GM should encourage players to tell brief in-character stories; the quality (or a Presence + Sway roll) determines the value.

\textbf{Mechanic:} When a player tells a story, they roll Presence + Sway (DV set by the listener's mood). On a success, the story is worth 1 Measure of value. On a critical success, it is worth 2. On a failure, the listener is offended — the party gains -1 Position on their next social roll in that circle. A player may also offer a \textbf{truth} (a secret, a confession, a memory) instead of a story; a truth is always worth 2 Measures but marks 1 Obligation to Ikasha.

\subsubsection*{Land Rent and Neutrality}
Tulkani circles are neutral ground. No lord may collect taxes or enforce laws inside a circle, regardless of whose territory it sits on. However, the lord may demand rent at the circle's edge, and the Tulkani will negotiate. A party inside a Tulkani circle is under the Caravan Judge's protection — but if the party commits a crime, the judge may expel them to the Hollow.

\textbf{Mechanic:} When a lord demands rent, the Caravan Judge mediates. The party may offer a story (one Measure), a promise (a future favour, represented by a Reciprocity Clock [4]), or a service (a quest). Refusing to pay means the lord may attack the circle — which breaks the neutrality and summons the Hollow.

\subsubsection*{Vow-Cords and Reciprocity}
Every Tulkani carries a vow-cord — a braided strand of horsehair and silk, with knots tied for every promise made and kept. A Tulkani who breaks a promise must cut the corresponding knot; the cut cord becomes a ghost that follows the oath-breaker until the promise is fulfilled. Outsiders may borrow a Tulkani vow-cord; the knots then record their promises. But the Tulkani see this not as debt but as \textbf{reciprocity}: what goes around comes around, and the circle always balances.

\textbf{Mechanic:} A character who accepts a vow-cord gains +1 die to any roll made to keep a promise related to the cord. However, if they break that promise, they must cut a knot — mark 1 Corruption and gain a Condition: \emph{Followed by a Ghost} (-1 die to stealth until the promise is fulfilled). A character may tie a new knot to make a new promise; each knot beyond the first increases the penalty for breaking. But a character who fulfills nine knots may ask the Caravan Judge for a \textbf{Circle Boon} — a permanent +1 die on all social rolls with Tulkani.

\subsubsection*{Local Patrons (Syncretic Names)}
\begin{itemize}
  \item \textbf{The Unwitnessed Weaver} — Inaea. *The knot that binds gently, the cord that remembers mercy.*
  \item \textbf{The Cord-Cutter} — Isoka. *The clean break, the shedding of false promises.*
  \item \textbf{The Dreamer Below} — Ikasha. *She who sleeps in the wagon, she who dreams the road.*
  \item \textbf{The Ninth Traveller} — The Traveler. *The road that has no end, the threshold that is always one step away.*
\end{itemize}

\subsection*{Lore Echoes (Cross-Expansion Hooks)}
\begin{itemize}
  \item \textbf{Tulkani Vow-Cords and the Sidhi Freedman's Seal:} Sidhi merchants on the Brass Coast accept Tulkani vow-cords as proof of credit. A well-knotted cord can buy a ship. A cut cord can sink one. The Sidhi know that a Tulkani knot is worth more than a signature.
  \item \textbf{Ikasha and the Veiled Sisters (Ashaan):} The Tulkani are the primary spreaders of Ikasha's cult outside Ashaan. A Breath-less agent may be disguised as a Tulkani story-seller; a Tulkani caravan may be a mobile safe house. The Sisters have tried to infiltrate the circles three times. Each time, the spy was found with a cord around her throat — tied in a knot she could not untie.
  \item \textbf{The Ninth Night and the Hollow (Eastern Realms):} The Tulkani's ninth-night rule is the only reason the Hollow has not overrun the eastern trade routes. Aelaerem hedge-keepers sometimes hire Tulkani guides to reset the Hollow's attention clocks. The payment is always a story — and the story is always about a door that was left open.
  \item \textbf{Story Debts and the Ukwe (Sekogo):} The Tulkani believe that the Ukwe is the first vow-cord tree. They leave offerings of knotted grass at its roots. A Tulkani judge may consult the Ukwe to verify a story's truth. The Ukwe speaks in rustles. A single rustle means true. Two mean false. Three mean the story is not finished.
\end{itemize}

\subsection*{Tulkani Superstitions (burn for +1 die once)}
\begin{itemize}
  \item \textbf{Bread to the Circle:} Crumble a crust at the centre of the wagon ring. Ignore the first \emph{Recognition} penalty. `[SUPERSTITION]`
  \item \textbf{Knot the Lantern:} Tie a single knot in the rope of a lantern. Cancel \emph{Lord's Rent Demand} or take –1 die on your next negotiation (your choice). `[SUPERSTITION]`
  \item \textbf{Name to the Ninth Post:} Whisper a dead judge's name into the central post. Gain +1 Effect when telling a story this scene, but mark \emph{Obligation} to Ikasha. `[SUPERSTITION]`
\end{itemize}

\subsection*{Thematic SB Spend Table}
\label{sec:tulkani-sb}

\paragraph{Minor Complications (1 SB)}
\begin{itemize}
\item \textbf{Exposure:} Your actions draw unwanted attention from \textbf{Tulkani judges or local lords}.
\item \textbf{Noise:} Sounds of your actions alert nearby \textbf{circle-watchers or Hollow scouts}.
\item \textbf{Trace:} Evidence of your passage marks your route for \textbf{vow-breakers or story-hunters}.
\item \textbf{Delay:} A brief but meaningful setback costs you \textbf{time or favourable camping ground}.
\item \textbf{Supply Strain:} Mark +1 segment on a relevant \textbf{resource clock}.
\end{itemize}

\paragraph{Moderate Setbacks (2 SB)}
\begin{itemize}
\item \textbf{Alarm Raised:} \textbf{Local Caravan Judge or lord's herald} becomes aware and begins responding.
\item \textbf{Position Lost:} You lose advantageous ground/cover/stealth due to \textbf{lantern failure or broken wheel}.
\item \textbf{Foe Appears:} A \textbf{debt collector or Hollow manifestation} arrives on scene.
\item \textbf{Gear Trouble:} A piece of equipment becomes \textbf{Compromised/Neglected}.
\item \textbf{Lock/Barrier:} A simple obstacle now requires a test to overcome.
\end{itemize}

\paragraph{Serious Trouble (3 SB)}
\begin{itemize}
\item \textbf{Reinforcements:} Additional \textbf{Tulkani riders, lord's men, or Hollow spirits} arrive.
\item \textbf{Key Gear Breaks:} A crucial tool/weapon becomes temporarily unusable.
\item \textbf{Major Twist:} The situation fundamentally changes—\textbf{knot cut/wagon broken/ninth night triggered}.
\item \textbf{Rail Tick:} Advance a relevant campaign/front clock by 1 segment.
\item \textbf{Condition Applied:} Mark \textbf{Fatigue 1/Harm 1/Condition} appropriate to fiction.
\end{itemize}

\paragraph{Major Turns (4+ SB)}
\begin{itemize}
\item \textbf{Trap Springs:} A prepared danger activates with full effect.
\item \textbf{Authority Arrival:} \textbf{High Caravan Judge, Ikasha's Knife, or Hollow's Debt-Walker} intervenes.
\item \textbf{Scene Shift:} The environment changes dramatically—\textbf{circle breaks/wagon burns/lanterns go out}.
\item \textbf{Patron Omen:} Divine/arcane forces take notice—\textbf{omen appears/blessing lost/curse manifests}.
\item \textbf{Narrative Pivot:} The story takes an unexpected turn that reframes objectives.
\end{itemize}

\paragraph{Region-Specific SB Options}
\begin{itemize}
\item \textbf{Tulkani (Knot Law):} Cords tighten without warning, a knot unties itself, a wheel stops at the ninth spoke.
\item \textbf{Tulkani (Story Debt):} A story you told returns to you changed, a listener laughs at a tragedy, the wind carries your words to someone who should not hear.
\item \textbf{Tulkani (Ninth Night):} A lantern flickers, a shadow stands at the circle's edge, a whisper says your name from the dark.
\end{itemize}

\subsection*{Pursuit Ladder (Near / Mid / Far)}
\begin{itemize}
\item \textbf{Advance/Withdraw (Wits + Stealth):} On a hit, shift one band. On a strong hit, also impose \emph{False Knot}.
\item \textbf{Cut the Cord (Presence + Sway):} On a hit, freeze range for one exchange; if a vow-bead is in your possession, +1 die.
\item \textbf{Weather the Hollow (Resolve + Spirit):} On a hit in \emph{Hollow Walks}, you pick who hears the whisper; on a miss, the GM does.
\end{itemize}

\subsection*{Boss Seeds (Knot \& Dream)}
\begin{itemize}
\item \textbf{The Cord-Cutter (Nine Broken Promises)} — \emph{Phases}: Vow-Cords Severed $\rightarrow$ Orchestrated Exile $\rightarrow$ Circle Coup. \emph{Cracks}: re-tie a cut cord, expose a false judge, survive a night outside the circle.
\item \textbf{The Dream-Eater (Ikasha's Shadow)} — \emph{Phases}: Dreamers Waking $\rightarrow$ Procession of Empty Wagons $\rightarrow$ Knot-Baptism of a Storyteller. \emph{Cracks}: relic dream-cord, survivor testimony, the ninth cord refusing to tighten.
\end{itemize}

\begin{tcolorbox}[colback=black!3,colframe=black!40!white,title={Patronage \& Power}]
Among the Tulkani, power flows through knots, stories, and the memory of the circle. A Caravan Judge's authority is recognised by every lord from the steppe to the coast — not because of armies, but because the Hollow respects the judge's cord. The Tulkani have no temples, no kings, no borders. They have the ninth night, and that is enough.

\textbf{For the GM:}  
Power in Tulkani society is measured in knots tied, stories told, and nights counted. Patronage often takes the form of vow-beads, dream-cords, or safe-haven oaths — each carrying the weight of a promise. To emphasize the itinerant mechanics:
\begin{itemize}
\item Require players to \textbf{negotiate land rent} with local lords, using stories or promises as currency. A failed negotiation may lead to the Hollow.
\item Use the \textbf{Ninth Night Clock} as a tension mechanic: every night in a Tulkani camp advances the clock. At 8, the Hollow manifests unless the party leaves.
\item Let \textbf{vow-cords} be physical props — actual knotted strings that players can hold, tie, or cut. A cut cord is a dramatic moment.
\item Emphasise that \textbf{storytelling is a skill}: reward creative players with mechanical bonuses, and let failed stories have consequences — a bored lord may raise rent, an offended spirit may curse.
\item Treat Tulkani society as a \textbf{reciprocity economy}, not a debt economy. What you give, you get back — but the circle decides when.
\end{itemize}
In Tulkani society, the knot remembers everything — and the ninth night forgets nothing.
\end{tcolorbox}

% Index entries
\index{Tulkani|see{Itinerant Courts}}
\index{Theme generator}
\index{Itinerant Courts generator}
\index{Places!Tulkani}
\index{People!Tulkani}
\index{Complications!Tulkani}
\index{Rewards!Tulkani}
\index{Ninth night}
\index{Vow-cords}
\index{Story as currency}
\index{Ikasha}
\index{Dream-cord}
\index{Caravan Judge}
\index{Wagon circle}
\index{Hollow}
\index{Tinkanar}
\index{Local Patrons!Tulkani}